\documentclass[12pt, a4paper, oneside]{report}
% ============================================================
%  Shared preamble — ViKi BS thesis (Innopolis University)
%  Compile with: xelatex (twice) — Cyrillic requires XeLaTeX
% ============================================================

\usepackage{fontspec}
\usepackage{polyglossia}
\setmainlanguage{english}
\setotherlanguage{russian}

\setmainfont{FreeSerif}[Ligatures=TeX]
\setsansfont{FreeSans}[Ligatures=TeX]
\setmonofont{DejaVu Sans Mono}[Scale=0.85]
\newfontfamily\cyrillicfont{FreeSerif}[Ligatures=TeX]
\newfontfamily\cyrillicfontsf{FreeSans}[Ligatures=TeX]
\newfontfamily\cyrillicfonttt{DejaVu Sans Mono}[Scale=0.85]

\usepackage{geometry}
\geometry{a4paper, top=25mm, bottom=25mm, left=30mm, right=15mm}

\usepackage{amsmath, amssymb, amsfonts, bm}
\usepackage{graphicx}
\usepackage{booktabs}
\usepackage{longtable}
\usepackage{array}
\usepackage{multirow}
\usepackage{caption}
\usepackage{subcaption}
\usepackage{enumitem}
\usepackage{xcolor}
\usepackage{setspace}
\usepackage{tikz}
\usetikzlibrary{arrows.meta, positioning, shapes.geometric, fit, backgrounds}
\usepackage[hidelinks]{hyperref}
\usepackage{tocloft}

\onehalfspacing
\setlength{\parindent}{1.25em}
\setlength{\parskip}{0.25em}

% --- IEEE-ish caption style -------------------------------
\captionsetup[figure]{labelfont=bf, labelsep=period, name=Fig., font=small}
\captionsetup[table]{labelfont=bf, labelsep=newline, justification=centering,
                     name=TABLE, font=small}
\renewcommand{\thetable}{\Roman{table}}

% --- placeholder macro ------------------------------------
% Every unresolved number / empirical claim is wrapped in \ph{}
% so it can be found with a single grep before submission.
\definecolor{phcol}{RGB}{176,0,32}
\newcommand{\ph}[1]{\textcolor{phcol}{\textbf{#1}}}
\newcommand{\TODO}[1]{\textcolor{phcol}{\small\textsf{[TODO: #1]}}}

% --- figure placeholder box -------------------------------
\newcommand{\figplaceholder}[2][0.62\textwidth]{%
  \begin{center}
  \fbox{\parbox[c][#2][c]{#1}{\centering\textcolor{phcol}{\textsf{\small
  [placeholder — insert figure]}}}}
  \end{center}}

% --- math notation ----------------------------------------
\newcommand{\R}{\mathbb{R}}
\newcommand{\SE}{\mathrm{SE}(3)}
\newcommand{\so}{\mathfrak{se}(3)}
\newcommand{\bq}{\bm{q}}
\newcommand{\bT}{\bm{T}}
\newcommand{\be}{\bm{e}}
\newcommand{\bx}{\bm{x}}
\newcommand{\bw}{\bm{w}}
\newcommand{\bJ}{\bm{J}}
\newcommand{\norm}[1]{\left\lVert #1 \right\rVert}
\DeclareMathOperator{\Log}{Log}
\DeclareMathOperator*{\argmin}{arg\,min}

% --- misc --------------------------------------------------
\newcommand{\viki}{ViKi}
\renewcommand{\cftsecleader}{\cftdotfill{\cftdotsep}}


\setmainlanguage{russian}
\setotherlanguage{english}

\begin{document}

\begin{titlepage}
\centering
\vspace*{10mm}
{\large УНИВЕРСИТЕТ ИННОПОЛИС}\\[2mm]
{\normalsize Бакалавриат по направлению «Компьютерные науки»}\\[35mm]

{\LARGE\bfseries Аннотация к выпускной\\[3mm]
квалификационной работе\par}
\vspace{10mm}

{\large\bfseries ViKi: кинематическая имитация на основе технического зрения\\[2mm]
для объектно-центричного ретаргетинга демонстраций\\[2mm]
манипуляции человека на роботизированные манипуляторы\par}
\vspace{25mm}

\begin{tabular}{rl}
\textit{Автор:} & \ph{Артём [Фамилия]}\\[4mm]
\textit{Научный руководитель:} & \ph{[ФИО] Корнаев}\\[4mm]
\textit{Группа:} & \ph{[группа]}\\
\end{tabular}

\vfill
{\normalsize Иннополис, Россия\\ \the\year}
\end{titlepage}

\tableofcontents
\cleardoublepage

% =========================================================
\chapter{Постановка задачи и актуальность}

\section{Предметная область}

Обучение роботизированных манипуляторов методом имитации (imitation learning)
стало практичным подходом к решению задач манипуляции благодаря появлению
визуомоторных политик, таких как Action Chunking Transformer (ACT) и Diffusion
Policy. Качество и обобщающая способность таких политик напрямую зависят от
объёма и качества обучающих демонстраций.

Доминирующим способом получения демонстраций сегодня является телеоперация:
человек-оператор проводит робота по траектории выполнения задачи, а состояния
суставов и видеопотоки записываются. Основное преимущество телеоперации в том,
что метки действий по построению исполнимы на целевом эмбодименте. Однако у
подхода есть структурные недостатки. Каждый эпизод требует времени оператора в
режиме реального времени; необходимое оборудование --- VR-контроллер,
экзоскелет или устройство с силовой обратной связью --- дорого и специфично для
конкретного эмбодимента; наконец, записанное движение определяется не столько
самой задачей, сколько задержками и эргономикой интерфейса управления.

Пассивная запись демонстраций, при которой человек выполняет задачу голыми
руками без какого-либо интерфейса, обращает этот компромисс. Сбор данных
становится дешёвым, быстрым и естественным: демонстратор успевает записать
десятки эпизодов за время, необходимое телеоператору на один. Утрачивается при
этом ровно то свойство, которое телеоперация гарантирует: отсутствуют метки
действий, а кинематика руки человека не соответствует кинематике
роботизированного манипулятора. Преодоление этого несоответствия и составляет
задачу ретаргетинга, которой посвящена настоящая работа.

\section{Проблема}

Пассивный захват удешевляет сбор данных, но порождает наблюдения, а не действия.
Преобразование наблюдаемой траектории кисти человека в траекторию суставов
робота, которая одновременно (i) воспроизводит значимое для задачи движение,
(ii) достаточно гладка для имитации политикой и (iii) физически исполнима на
целевом манипуляторе, представляет собой недоопределённую задачу, осложнённую
шумом восприятия.

Доминируют два режима отказа.

\textbf{Усиление шума через кинематическую цепь.} Монокулярные или нейросетевые
оценки глубины ключевых точек кисти зашумлены на сантиметровом масштабе.
Обратная задача кинематики нелинейна, поэтому ограниченное возмущение в
пространстве задачи может отобразиться в неограниченное отклонение в
пространстве суставов вблизи сингулярностей. Широко применяемый ответ ---
низкочастотная фильтрация наблюдаемой траектории до решения обратной кинематики
с последующей фильтрацией траектории суставов --- проблему не решает.
Предварительная фильтрация применяется в пространстве задачи, тогда как
гладкость требуется в пространстве суставов, и вдобавок она удаляет истинное
высокочастотное содержание задачи вместе с шумом. Именно эта схема
«мусор на входе --- мусор на выходе» и мотивирует настоящую работу.

\textbf{Представленческая зависимость.} Траектория, записанная как абсолютная
последовательность положений в системе координат камеры или мира, привязана к
конкретной пространственной раскладке сеанса записи. Если объект переместить,
траектория окажется неверной; при смене эмбодимента --- неверной иначе.
Обобщается отношение между кистью и манипулируемым объектом, а не абсолютный
путь кисти в пространстве.

\section{Научный пробел}

Оба режима отказа имеют известные решения, но эти решения приняты в разных,
не пересекающихся сообществах.

Оптимизационный ретаргетинг, в котором точность соответствия данным, временная
гладкость и ограничения на суставы формулируются как единая задача условной
оптимизации по конфигурации робота, стал стандартом в системах
\emph{онлайн-телеоперации}. Каноническую форму задали AnyTeleop и связанная с
ним библиотека \texttt{dex-retargeting}; на них непосредственно опираются
Open-TeleVision, Bunny-VisionPro и ACE; DexFlow расширяет штраф гладкости до
второго порядка. Однако во всех случаях человек остаётся в контуре управления,
оптимизация выполняется на интерактивных частотах, и датасет не формируется.

Напротив, \emph{офлайн}-пайплайны, генерирующие датасеты из пассивного видео
человека, принимают объектно-центричные представления, но не единую
оптимизационную постановку. ORION планирует по разреженным ключевым кадрам
объекта и вовсе не выполняет покадровый ретаргетинг кисти. OKAMI восстанавливает
параметрическую модель тела и кисти и выполняет ретаргетинг в две развязанные
стадии --- варпинг референсной траектории под новое положение объекта с
последующим решением обратной кинематики, --- что структурно совпадает с
раскритикованной выше последовательной схемой.

Целевой анализ предшествующих работ, проведённый в рамках настоящего
исследования, не выявил опубликованной системы, объединяющей одновременно:
пассивный безмаркерный захват; метрически откалиброванный мультикамерный RGB-D
риг; объектно-центричное плотное траекторное представление; ретаргетинг как
единый функционал стоимости; экспорт датасета, готового к обучению политик; и
физическое воспроизведение как этап валидации до экспорта. Ближайшие системы
удовлетворяют части критериев и не удовлетворяют остальным. Таким образом,
закрываемый пробел состоит не в изобретении нового оператора ретаргетинга, а в
переносе устоявшейся формулировки в режим, где она не применялась, вместе с
необходимой этому режиму сенсорной и валидационной инфраструктурой.

\section{Исследовательские вопросы и гипотезы}

Работа направляется одним основным исследовательским вопросом и тремя
подвопросами.

\vspace{2mm}
\noindent\textbf{ИВ.} \emph{Каким образом демонстрации манипуляции, захваченные
пассивно откалиброванным мультикамерным RGB-D ригом, могут быть подвергнуты
ретаргетингу на роботизированный манипулятор так, чтобы получить датасеты,
пригодные для обучения визуомоторных имитационных политик?}

\begin{enumerate}[label=\textbf{ИВ\arabic*.}, leftmargin=2.4em]
  \item Как постановка ретаргетинга в виде единого взвешенного функционала
        стоимости, объединяющего член соответствия данным со взвешиванием по
        достоверности с штрафами на скорость и ускорение, влияет на точность и
        гладкость результирующих траекторий суставов по сравнению с
        последовательной схемой smooth--IK--smooth?
  \item Какая точность оценки позы запястья достижима при
        взвешенном по достоверности слиянии нескольких откалиброванных RGB-D
        видов, привязанных к общей системе координат рабочей области по
        маркерной доске ChArUco, и как остаточная ошибка декомпозируется по
        независимым источникам?
  \item В какой степени воспроизведение синтезированных траекторий на
        физическом манипуляторе до экспорта датасета повышает исполнимость
        экспортируемых данных?
\end{enumerate}

\noindent\textbf{Гипотезы.}

\textbf{Г1 (поведение).} Ретаргетинг, поставленный как единый взвешенный
функционал стоимости, даёт траектории суставов с меньшей ошибкой отслеживания в
пространстве задачи и большей гладкостью, чем последовательная схема
smooth--IK--smooth, на одних и тех же записанных демонстрациях. Независимая
переменная --- формулировка ретаргетинга; зависимые --- RMSE отслеживания позы
запястья и спектральная длина дуги траектории суставов; популяция ---
записанные демонстрации ограниченного класса задач; отношение --- доминирование
единой формулировки по обеим мерам.

\textbf{Г2 (поведение).} Взвешивание члена соответствия данным по покадровой
достоверности восприятия снижает ошибку отслеживания относительно равномерного
взвешивания, причём снижение больше на кадрах с частичной окклюзией ключевых
точек кисти.

\textbf{Г3 (надёжность).} Отбор траекторий по результатам физического
воспроизведения до экспорта повышает долю экспортированных эпизодов, успешно
исполняемых на манипуляторе.

Вопросы проверены по критериям FINER. Они \emph{выполнимы} --- оборудование,
решатель и протокол оценки доступны, класс задач намеренно ограничен;
\emph{интересны} и \emph{релевантны} сообществу робототехнического обучения,
поскольку стоимость демонстраций является признанным узким местом;
\emph{новы} в установленном выше смысле переноса режима, а не нового оператора;
и \emph{этичны}, поскольку исследование не вовлекает испытуемых помимо автора и
не сохраняет персональных данных.

% =========================================================
\chapter{Обзор литературы}

\section{Представление кисти человека}

Значительная часть работ восстанавливает полную кинематику демонстратора. OKAMI
сочетает SMPL-H для тела с MANO для кисти и выполняет ретаргетинг раздельно.
DexCanvas строит датасет соответствий человек--робот, подгоняя MANO к данным
двадцати двух инфракрасных камер захвата движения и двух RGB-D сенсоров. DexMV
аналогично восстанавливает полную модель MANO вместе с позой объекта.

Полная реконструкция даёт богатое представление и необходима, когда целевым
эмбодиментом является многопальцевая антропоморфная кисть. Одновременно она
вносит тяжёлый перцептивный стек и дополнительный, плохо охарактеризованный
источник ошибки --- саму подгонку меша. Для параллельного схвата большая часть
этого представленческого богатства далее отбрасывается.

Альтернатива, принятая в настоящей работе, представляет кисть как позу запястья
в $\SE$ вместе с бинарным состоянием захвата. Это не уступка ради
вычислительной простоты, а представление, используемое магистральным
направлением современных систем телеоперации. Open-TeleVision отображает позу
запястья оператора на схват манипулятора через замкнутую обратную кинематику и
кодирует параллельный захват единственным вектором большой--указательный палец;
Bunny-VisionPro и ACE принимают ту же декомпозицию, причём ACE дополнительно
демонстрирует её переносимость между антропоморфными кистями, параллельными
схватами и четвероногими роботами.

Методологическую поддержку ослаблению веса запястья даёт абляционное
исследование Xin и соавторов, которые собирают всеобъемлющий функционал
ретаргетинга --- глобальная поза кисти, общая форма кисти, относительные
положения кончиков пальцев, ориентации кончиков пальцев --- и удаляют члены по
одному. Они обнаруживают, что точное отслеживание глобальной позы запястья
часто избыточно и способно ухудшать точность по кончикам пальцев, и заменяют
позиционную цель по запястью целью по кончику большого пальца, регуляризуя
ориентацию запястья лишь слабо. Для представления «запястье плюс захват» отсюда
следует, что цель по запястью должна входить в оптимизацию как взвешенная мягкая
задача, а не как жёсткое ограничение.

\section{Формулировки ретаргетинга}

Преобладающая современная формулировка ставит ретаргетинг как задачу условных
нелинейных наименьших квадратов или квадратичного программирования по
конфигурации робота на каждом шаге, минимизируя взвешенную сумму члена
соответствия данным и члена временной гладкости при ограничениях на пределы
суставов. AnyTeleop даёт каноническую реализацию: целевая функция суммирует по
ключевым точкам квадрат невязки между масштабированным вектором ключевой точки
человека и соответствующим вектором прямой кинематики робота, плюс штраф первого
порядка на изменение конфигурации между соседними кадрами.

Два уточнения непосредственно релевантны настоящей работе. Первое --- покадровое
взвешивание членов: в формулировке Open-TeleVision невязка каждой ключевой точки
масштабируется переключающим весом, зависящим от расстояния и возрастающим при
обнаружении щипкового захвата. Механизм в точности аналогичен взвешиванию по
покадровой достоверности восприятия. Второе --- гладкость второго порядка:
DexFlow дополняет стандартный функционал динамическим весом сглаживания и
регуляризацией гессиана для обеспечения $C^2$-непрерывности, штрафуя кривизну, а
не только скорость. Робастификация функцией потерь Хьюбера на члене данных даёт
механизм отбраковки выбросов восприятия внутри оптимизатора, а не сглаживания их
заранее.

Последовательная схема --- фильтрация наблюдаемой траектории в пространстве
задачи, покадровое решение обратной кинематики, фильтрация выхода в пространстве
суставов --- остаётся распространённой в офлайн-пайплайнах. Её дефект
структурный, а не эмпирический: фильтрация применяется в пространстве задачи,
тогда как требование гладкости относится к траектории суставов, и, поскольку
прямая кинематика нелинейна, гладкий декартов вход не влечёт гладкого выхода в
пространстве суставов, особенно вблизи сингулярностей. В единой функциональной
постановке штраф гладкости действует непосредственно на конфигурацию, поэтому
гладкость в пространстве суставов достигается по построению.

\section{Представление задачи}

Объектно-центричные представления доминируют в недавних работах по обобщению
между окружениями и эмбодиментами. ORION представляет задачу как трёхмерную
траекторию объекта относительно начального и конечного положений объекта,
утверждая, что это обобщается на смену визуального фона, сдвиг камеры,
пространственную раскладку и новые экземпляры объектов, и что полная
реконструкция тела демонстратора не требуется. SPOT кодирует задачу как
$\SE$-траекторию позы объекта относительно цели именно с целью развязать
эмбодимент-специфичные действия от сенсорного входа, что позволяет обучаться на
демонстрациях без действий. Подходы на основе потока приводят эквивалентный
аргумент, ослабляя при этом предположение о жёсткости объекта.

Компромисс состоит в том, что объектно-центричные представления требуют
устойчивой оценки позы или сегментации объекта, а варианты на основе траектории
позы предполагают жёсткость, тогда как траектории в мировой системе координат
проще, но привязаны к раскладке сеанса записи. Хранение обоих представлений
переносит выбор на потребителя данных при пренебрежимой стоимости хранения.

\section{Калибровка, слияние и синхронизация мультикамерного RGB-D рига}

Доска ChArUco сочетает субпиксельную точность углов шахматной доски с
устойчивой идентификацией и толерантностью к окклюзии, свойственными маркерам
ArUco. Доска $6\times8$ даёт до тридцати пяти углов на вид против четырёх у
одиночного маркера AprilTag и остаётся детектируемой при частичной окклюзии и на
границе кадра. Типичные ошибки репроекции для калибровки ChArUco и ArUco лежат
ниже половины пикселя. Маркеры AprilTag проще детектируются на больших
дистанциях и при наклонном падении и достаточны для оценки только внешних
параметров, но плохо подходят для калибровки внутренних параметров из-за малого
числа углов.

Romeo и соавторы дают наиболее чистое прямое сравнение процедур калибровки для
мульти-Kinect ригов, сообщая, что трёхмерные процедуры превосходят двумерные
подходы на основе шахматной доски: средние расстояния между соответствующими
точками облаков составили $21{,}4$~мм при цветовом разрешении и $9{,}9$~мм при
инфракрасном для статичной сцены и, соответственно, $20{,}9$~мм и $7{,}4$~мм для
динамичной. Bu и Lee сообщают --- что важно для геометрии рига, --- что
уточнение методом итеративной ближайшей точки не сходится для почти
противоположных ракурсов камер из-за недостаточного перекрытия облаков точек.

Спецификации одиночных сенсоров ограничивают достижимую точность снизу. Azure
Kinect DK специфицирован стандартным отклонением случайной ошибки не более
$17$~мм и типичной систематической ошибкой ниже $11$~мм плюс $0{,}1$~\%
дистанции; независимые оценки подтверждают эти значения. Практическое следствие
состоит в том, что между откалиброванными видами следует ожидать согласования
сантиметрового, а не миллиметрового уровня.

Наивное усреднение перекрывающихся измерений глубины уступает взвешенному по
достоверности слиянию. Взвешивание по дистанции и по углу между нормалью
поверхности и направлением наблюдения --- так, что поверхности, наблюдаемые
близко к перпендикуляру, вносят больший вклад, чем наблюдаемые под скользящим
углом, --- является устоявшейся практикой.

Ловушка для гетерогенного рига состоит в том, что аппаратные часы устройств
разных семейств взаимно несопоставимы, поскольку каждое считает собственные
циклы от собственной эпохи. Документированная практика для таких ригов ---
маркировать каждый поступающий кадр единым host-side монотонным таймстампом и
выравнивать потоки по ближайшему соседу в этом общем домене, принимая точность
миллисекундного масштаба. Поскольку триггеры Azure Kinect и RealSense взаимно
несовместимы, это единственный доступный механизм выравнивания двух семейств.

\section{Позиционирование предлагаемой системы}

Сравнение проводится по шести критериям, совместно определяющим предлагаемую
архитектуру: (К1) пассивный безмаркерный захват без телеоперационного
оборудования; (К2) мультикамерное метрически откалиброванное RGB-D зондирование;
(К3) объектно-центричное плотное траекторное представление; (К4) ретаргетинг как
единый взвешенный по достоверности функционал стоимости; (К5) экспорт датасета,
готового к имитационному обучению; (К6) физическое воспроизведение для валидации
датасета до экспорта.

Ни одна из обнаруженных систем не удовлетворяет всем шести критериям
одновременно. HOWTransfer ближе всего по последовательности операций,
удовлетворяя К1 и К6, однако является кисте-центричной, а не
объектно-центричной, использует откалиброванную стереопару RGB вместо
метрического RGB-D и не применяет единый функционал стоимости. Vi-RoIL ближе
всего по представлению, удовлетворяя К3, но восстанавливает метрический масштаб
из монокулярного вида маркерного куба и оценивается только в симуляции.
Пересечение литературы по онлайн-оптимизационному ретаргетингу с литературой по
офлайн-генерации объектно-центричных датасетов остаётся незанятым, и именно на
него нацелена настоящая работа.

% =========================================================
\chapter{Предлагаемое решение}

\section{Общая структура пайплайна}

ViKi --- офлайн-пайплайн: демонстрации записываются, обрабатываются и
экспортируются раздельными стадиями, без требования выполнения какой-либо стадии
на интерактивных частотах. Это осознанный выбор. Ослабление ограничения
реального времени, формирующего архитектуру систем телеоперации, допускает
глобальную оптимизацию по целым траекториям, итеративное уточнение и стадию
валидации, исполняемую на физическом оборудовании, --- ни одно из этих средств
недоступно онлайн-системе.

Пайплайн состоит из шести стадий: захват синхронизированных RGB-D потоков;
восприятие с детекцией ключевых точек кисти, подъёмом в три измерения по данным
глубины и слиянием видов со взвешиванием по достоверности; представление
траектории запястья относительно объекта и относительно якоря рабочей области;
ретаргетинг путём минимизации единого взвешенного функционала при кинематических
и коллизионных ограничениях; валидация воспроизведением на физическом
манипуляторе; экспорт в LeRobot-совместимом формате.

\section{Калибровка и синхронизация}

Внутренние параметры всех цветовых и глубинных потоков оцениваются по
наблюдениям доски ChArUco. Внешняя калибровка следует принципу якоря рабочей
области: вместо выражения камер друг относительно друга каждая камера
калибруется относительно доски ChArUco, жёстко закреплённой в рабочей области и
задающей общую мировую систему координат. Привязка к рабочей области, а не к
опорной камере, имеет два следствия: добавление или перемещение камеры не
аннулирует калибровку остальных, а база манипулятора может быть
зарегистрирована в ту же систему координат без отдельной цепочки
камера--камера.

Две камеры Azure Kinect синхронизируются аппаратно в конфигурации
ведущий/ведомый, причём ведомые устройства запускаются раньше ведущего, а
экспозиции глубины разносятся во избежание взаимной засветки лазеров. Камера
RealSense не может разделять этот триггер. Межсемейственное выравнивание
поэтому использует единый host-side монотонный таймстамп, присваиваемый каждому
кадру при поступлении, с последующим сопоставлением потоков по ближайшему
соседу. Подчеркнём два момента: аппаратные таймстампы устройств не используются
для межустройственной арифметики, а монотонные часы применяются вместо часов
реального времени, поскольку последние могут быть скачкообразно скорректированы
службой синхронизации времени, что испортило бы вычисление интервалов посреди
записи.

\section{Слияние со взвешиванием по достоверности}

Оценки от нескольких видов объединяются взвешенным по достоверности усреднением,
а не наивным средним. Вес каждой ключевой точки в каждом виде образуется
произведением четырёх факторов: оценки видимости детектора; достоверности или
флага валидности глубины сенсора в соответствующем пикселе; обратного квадрата
дальности; и косинуса угла падения, отсечённого снизу нулём. Произведённая форма
выбрана так, чтобы нуль в любом факторе --- недетектированная точка, невалидный
пиксель глубины, скользящий ракурс --- полностью устранял вклад данного вида, а
не просто ослаблял его.

\section{Ретаргетинг как единый функционал стоимости}

Пусть $\bq_t \in \R^{n}$ --- конфигурация манипулятора в момент $t$, а
$f(\bq_t) \in \SE$ --- прямая кинематика схвата. Невязка в пространстве задачи
относительно желаемой позы есть логарифмическое отображение ошибки позы,
$\be_t(\bq_t) = \Log\big((\bT^{R,\mathrm{des}}_{E,t})^{-1} f(\bq_t)\big)^{\vee}
\in \R^{6}$. Ретаргетинг ставится как единая условная задача

\begin{equation}
  \min_{\bq_{1:T}} \sum_{t=1}^{T} \Big[
    \omega_t \rho_\delta\!\left(\norm{\be_t(\bq_t)}_{\bm{\Sigma}}\right)
    + \lambda_v \norm{\bq_t - \bq_{t-1}}^{2}
    + \lambda_a \norm{\bq_t - 2\bq_{t-1} + \bq_{t-2}}^{2}
    + \lambda_r \norm{\bq_t - \bq_{\mathrm{ref}}}^{2}
  \Big]
\end{equation}

при ограничениях на пределы суставов, пределы скоростей и барьерных функциях
коллизий, где $\omega_t$ --- покадровый вес достоверности восприятия,
$\rho_\delta$ --- функция потерь Хьюбера, а $\lambda_v$, $\lambda_a$,
$\lambda_r$ --- веса штрафов на скорость, ускорение и отклонение от опорной
позы.

Покадровый вес агрегирует достоверности ключевых точек, породивших оценку
запястья в момент $t$. Конструкция является офлайн-аналогом переключающих весов,
зависящих от расстояния, применяемых в онлайн-ретаргетинге: кадры, в которых
кисть частично окклюдирована или плохо наблюдаема, оказывают пропорционально
меньшее влияние на решение, а члены гладкости проводят траекторию через них
вместо того, чтобы член данных стягивал её к ложному наблюдению.

Функция потерь Хьюбера ограничивает влияние грубых выбросов --- ошибочно
детектированной кисти, провала глубины, переживших взвешивание по достоверности,
--- на уровне самого члена данных, а не позволяет стадии фильтрации размазать
выброс по соседним кадрам. Именно этот механизм заменяет предварительное
сглаживание.

Ключевое структурное отличие от последовательной альтернативы состоит в том, что
штрафы гладкости действуют на конфигурацию, то есть в том пространстве, где
гладкость и требуется, и балансируются против соответствия данным самим
решателем, а не применяются вслепую заранее. Компромисс становится явным: на
кадрах высокой достоверности вес велик и решение отслеживает наблюдение; на
кадрах низкой достоверности доминируют члены гладкости и решение интерполирует.

Задача решается средствами Pink на Pinocchio. Каждый член отображается на
взвешенную задачу квадратичного программирования, тогда как ограничения входят
как неравенства, а избегание коллизий выражается через барьерные функции
управления. Демпфирование Левенберга--Марквардта стабилизирует решение вблизи
сингулярностей. Поскольку постановка офлайновая, задача решается по всей
траектории, а не покадрово, поэтому член ускорения корректно определён на каждом
внутреннем шаге, а решение не стеснено требованием причинности, с которым
сталкивается онлайн-решатель.

Динамические фильтры, налагающие модель движения, --- такие как фильтр Калмана
или расширенный фильтр Калмана, --- намеренно не применяются: они вносят
априорное предположение о динамике, которое затем переоценивается оптимизатором,
а взаимодействие двух априорных предположений трудно охарактеризовать.

\section{Валидация воспроизведением}

Синтезированные траектории исполняются на физическом манипуляторе до
финализации датасета. Стадия служит трём целям. Во-первых, отбраковке
нереализуемости: траектория может удовлетворять кинематическим ограничениям в
модели и всё же отказать на оборудовании из-за пределов отслеживания
контроллера, немоделированных коллизий или граничных условий рабочей области.
Во-вторых, закрытию разрыва эмбодиментов у источника: проприоцептивные
состояния, записанные при воспроизведении, суть те, которых робот действительно
достиг, а не те, что предсказал оптимизатор. Экспорт записанных, а не
синтезированных состояний означает, что датасет содержит честные пары
наблюдение--действие с целевого эмбодимента, --- ровно то свойство, которым
телеоперационные датасеты обладают по построению, а наивно выведенные из видео
--- нет. В-третьих, формированию сигнала курации: эпизоды ранжируются по
невязке отслеживания при воспроизведении.

% =========================================================
\chapter{Реализация, результаты и обсуждение}

\section{Реализация}

Система реализована на Python без слоя робототехнического middleware. Отказ от
ROS был осознанным: пайплайн представляет собой систему пакетной обработки
данных, а не распределённую систему управления реального времени, и граф узлов,
сериализация сообщений и инструментарий сборки, предоставляемые middleware,
добавили бы эксплуатационную сложность, не отвечая ни на одно требование
офлайн-постановки. Взаимодействие с манипулятором ведётся напрямую через
\texttt{ur-rtde}.

Реализация разделена на модули захвата, калибровки, восприятия, представления,
ретаргетинга, воспроизведения, экспорта и интерфейса. Разделение следует стадиям
пайплайна, причём контракт систем координат --- каждая пространственная величина
несёт явный идентификатор системы координат --- обеспечивается на границах
модулей, так что ошибки систем координат проявляются как ошибки типов, а не как
безмолвные численные ошибки. Система распространяется под лицензией MIT.

\section{Экспериментальная установка}

Эксперименты выполнены на манипуляторе UR3 с параллельным схватом, двух камерах
Azure Kinect DK в аппаратной синхронизации ведущий/ведомый и одной Intel
RealSense D435i, наблюдающих настольную рабочую область, привязанную доской
ChArUco. Расположение камер близко к ортогональному ради сохранения перекрытия
полей зрения.

Оценка проводилась на намеренно ограниченном классе задач, выбранном так, чтобы
выполнялись квазистатическое предположение и предположение о жёсткости объекта, а
успех измерялся однозначно: захват и перенос объекта в отмеченную целевую
позицию; то же с объектом, инициализированным в позе, не встречавшейся при
демонстрации, для проверки объектно-относительного представления; переливание,
проверяющее точность по ориентации, а не только по положению.

Предлагаемая формулировка сравнивается с последовательной базовой схемой, а
также с абляциями, каждая из которых отличается от полной системы ровно одним
компонентом: без взвешивания по достоверности; без члена ускорения; без
робастификации Хьюбера; с представлением в мировой системе координат вместо
объектно-относительного.

\section{Метрики}

Метрики зафиксированы до финальных прогонов, чтобы критерии успеха не
подбирались после наблюдения результатов. Точность ретаргетинга измеряется
среднеквадратичной ошибкой отслеживания в пространстве задачи, раздельно по
трансляции и вращению. Гладкость измеряется спектральной длиной дуги профиля
скоростей суставов --- безразмерной мерой, инвариантной к амплитуде и
длительности, --- вместе со средним модулем рывка. Реализуемость измеряется
долей синтезированных эпизодов, исполняемых на оборудовании без отказа
контроллера, нарушения пределов суставов или коллизии. Успешность нижележащих
политик измеряется для ACT и Diffusion Policy, обученных на экспортированном
датасете. Точность калибровки измеряется ошибкой репроекции, межкамерной
ошибкой регистрации и остаточным рассогласованием таймстампов между семействами
устройств.

\section{Результаты}

\ph{[Раздел заполняется по завершении экспериментов.]}

Предварительно ожидается, что предлагаемая формулировка снизит ошибку
отслеживания позы запястья на \ph{XX~\%} и повысит гладкость в пространстве
суставов на \ph{XX~\%} относительно последовательной базовой схемы, что отбор по
результатам воспроизведения отклонит \ph{XX~\%} эпизодов как нереализуемые, а
политики, обученные на полученном датасете, достигнут успешности \ph{XX~\%}.

\section{Декомпозиция ошибки}

Методологическое обязательство настоящей работы состоит в том, что остаточная
ошибка декомпозируется, а не сообщается единой агрегированной величиной. Три
вклада независимы по происхождению и требуют разных мер; их смешение сделало бы
агрегат неинформативным относительно того, что именно следует улучшать.

\emph{Шум детекции ключевых точек} --- ошибка в положениях ключевых точек кисти
на плоскости изображения, распространяемая через подъём по глубине. Этот вклад
снижается взвешиванием по достоверности и мультивидовым слиянием и не снижается
улучшением калибровки.

\emph{Ошибка калибровки} --- ошибка в преобразованиях камера--рабочая область и
робот--рабочая область. Этот вклад систематичен в пределах сеанса, проявляется
как постоянное смещение, а не как шум, и не снижается мультивидовым слиянием,
поскольку все виды его наследуют.

\emph{Представленческая зависимость} --- ошибка, вносимая выбором системы
отсчёта, в частности чувствительность траекторий в мировой системе координат к
смещению объекта. Этот вклад равен нулю, когда объект находится там же, где при
демонстрации, и растёт со смещением.

\section{Ограничения}

Бинарное состояние захвата не способно представлять внутрикистевую манипуляцию,
перехват или захваты с модуляцией усилия. Объектно-относительное преобразование
не определено для деформируемых объектов, что исключает манипуляцию тканью,
кабелем и сыпучими средами. Объектно-относительное представление наследует
ошибку оценщика позы объекта и при его отказе --- на бестекстурных, симметричных
или окклюдированных объектах --- деградирует хуже той альтернативы в мировой
системе координат, которую призвано улучшить. Силы и контактная динамика не
оцениваются, поэтому успех на задачах с компонентой регулирования усилия не
заявляется. Оценка использует одного демонстратора и ограниченный класс задач,
поэтому обобщение по демонстраторам, морфологиям кисти и семействам задач не
установлено. Наконец, согласование между камерами на сантиметровом уровне
ограничивает достижимую точность независимо от формулировки ретаргетинга.

% =========================================================
\chapter{Выводы}

Работа рассмотрела вопрос о том, каким образом демонстрации манипуляции,
записанные пассивно, могут быть преобразованы в датасеты, пригодные для обучения
визуомоторных имитационных политик. Показано, что два устоявшихся подхода к этой
задаче по отдельности неполны: системы онлайн-телеоперации сошлись к
методологически корректной формулировке ретаргетинга --- единому взвешенному
функционалу стоимости, --- но не порождают датасета, тогда как офлайн-пайплайны
порождают датасеты, но выполняют ретаргетинг либо вовсе не выполняя его, через
разреженные ключевые кадры объекта, либо через последовательную схему
фильтрация--решение--фильтрация, сглаживание в которой применяется не в том
пространстве, где требуется гарантируемое ею свойство.

Предложена система ViKi, занимающая пересечение этих подходов. Она захватывает
демонстрации метрически откалиброванным мультикамерным RGB-D ригом,
привязанным к системе координат рабочей области по доске ChArUco; представляет
кисть минимально --- позой запястья в $\SE$ с бинарным состоянием захвата;
хранит каждую демонстрацию и относительно манипулируемого объекта, и
относительно якоря рабочей области; выполняет ретаргетинг минимизацией единого
взвешенного функционала, в котором член данных масштабируется покадровой
достоверностью восприятия и робастифицируется функцией Хьюбера, а штрафы на
скорость и ускорение действуют непосредственно в пространстве суставов;
валидирует результат исполнением на физическом манипуляторе до экспорта; и
записывает LeRobot-совместимый датасет, содержащий проприоцепцию, которой робот
действительно достиг, а не состояния, предсказанные оптимизатором.

Следует отдельно указать, что \emph{не} заявляется. Функционал стоимости не нов
как оператор; его компоненты присутствуют в предшествующих работах, и заявление
об алгоритмической новизне было бы опровергнуто этим prior art.
Объектно-центричное представление, привязка по ChArUco, взвешенное по
достоверности слияние глубины и синхронизация по монотонным часам хоста ---
устоявшиеся практики. Вклад состоит в конъюнкции и переносе режима, которые
проведённый анализ предшествующих работ обнаружил незанятыми.

\vspace{4mm}
\noindent\ph{\small Примечание: величины и утверждения, выделенные этим цветом,
подлежат заполнению или проверке по завершении экспериментальной части. Полный
список источников приведён в основном тексте работы на английском языке.}

\end{document}
